% Part: propositional-logic
% Chapter: syntax-and-semantics
% Section: formulas

\documentclass[../../../include/open-logic-section]{subfiles}

\begin{document}

\olfileid{pl}{syn}{fml}

\olsection{विधानीय \usetoken{P}{formula}}

विधानीय तर्कशास्त्रातील !!^{formula}s ही
\emph{!!{propositional
    variable}s}\iftag{prvFalse}{\iftag{prvTrue}{,}{ आणि} विधानीय
  स्थिर~$\lfalse$}{}\iftag{prvTrue}{\iftag{prvFalse}{
    आणि}{} विधानीय स्थिर~$\ltrue$}{} यांपासून \emph{तार्किक
  संयोजके} वापरून तयार होतात.

\begin{enumerate}
\item !!^a{denumerable} संच~$\PVar$: त्यातील !!{propositional variable}s
  $\Obj p_0$, $\Obj p_1$, \dots
\tagitem{prvFalse}{!!{falsity} साठीचे विधानीय स्थिर~$\lfalse$.}{}
\tagitem{prvTrue}{!!{truth} साठीचे विधानीय स्थिर~$\ltrue$.}{}
\item तार्किक संयोजके:
  \startycommalist
  \iftag{prvNot}{\ycomma $\lnot$ (नकरण)}{}%
  \iftag{prvAnd}{\ycomma $\land$ (संधियोग)}{}%
  \iftag{prvOr}{\ycomma $\lor$ (विकल्पयोग)}{}%
  \iftag{prvIf}{\ycomma $\lif$ (!!{conditional})}{}%
  \iftag{prvIff}{\ycomma $\liff$ (!!{biconditional})}{}%
\item विरामचिन्हे: (, ), आणि स्वल्पविराम.
\end{enumerate}

विधानीय तर्कशास्त्राची ही भाषा आपण $\Lang L_0$ ने दर्शवतो.

\iftag{defNot,defOr,defAnd,defIf,defIff,defTrue,defFalse,defEx,defAll}{%
वर दिलेल्या आदिम संयोजकांबरोबर आपण पुढील \emph{परिभाषित} चिन्हेही
वापरतो:
\startycommalist
  \iftag{defNot}{\ycomma $\lnot$ (नकरण)}{}%
  \iftag{defAnd}{\ycomma $\land$ (संधियोग)}{}%
  \iftag{defOr}{\ycomma $\lor$ (विकल्पयोग)}{}%
  \iftag{defIf}{\ycomma $\lif$ (!!{conditional})}{}%
  \iftag{defIff}{\ycomma $\liff$ (!!{biconditional})}{}%
  \iftag{defFalse}{\ycomma $\lfalse$ (!!{falsity})}{}%
  \iftag{defTrue}{\ycomma $\ltrue$ (!!{truth})}}{}.

\begin{tagblock}{defNot,defOr,defAnd,defIf,defIff,defTrue,defFalse,defEx,defAll}
\begin{explain}
परिभाषित चिन्ह हे अधिकृतपणे भाषेचा भाग नसते; त्याऐवजी, अनौपचारिक
संक्षेप म्हणून त्याची ओळख करून दिली जाते. केवळ आदिम चिन्हे वापरल्यास
बरीच लांब होणारी सूत्रे त्यामुळे संक्षेपाने लिहिता येतात. हा लाभ
उघडच आहे. पण त्याहून मोठा लाभ असा की सिद्धता लहान होतात. एखादे चिन्ह
आदिम असेल, तर सिद्धतांमध्ये त्याचा स्वतंत्रपणे विचार करावा लागतो.
त्यामुळे आदिम चिन्हे जितकी अधिक, तितक्या आपल्या सिद्धता लांब होतात.
\end{explain}
\end{tagblock}

% Alternate symbols

\begin{intro}
वर आपण वापरलेल्या संज्ञांपेक्षा आणि चिन्हांपेक्षा वेगळ्या संज्ञा व
चिन्हे तुम्हाला परिचित असू शकतात. तर्कशास्त्राच्या ग्रंथांत (आणि
अध्यापनात) ``नकरण'' यासाठी सामान्यतः $\sim$, $\neg$, आणि~!\ यांपैकी
एखादे, तर ``संधियोग'' यासाठी $\wedge$, $\cdot$, आणि $\&$ यांपैकी
एखादे चिन्ह वापरतात. ``सशर्त (संकेतार्थक)'' किंवा ``अभिव्यंजन'' यांसाठी
$\rightarrow$, $\Rightarrow$, आणि $\supset$ ही चिन्हे सामान्यतः
वापरली जातात.
\iftag{prvIff,defIff}{``द्विसशर्त'', ``द्वि-अभिव्यंजन'' किंवा
  ``(वास्तविक) सममूल्यता'' यांसाठी $\leftrightarrow$,
  $\Leftrightarrow$, आणि $\equiv$ ही चिन्हे वापरतात.}{}
\iftag{prvFalse,defFalse}{$\lfalse$ या चिन्हाला वेगवेगळ्या ठिकाणी
  ``असत्यता'', ``फाल्सम'', ``विसंगती'' किंवा ``तळ'' म्हणतात.}{}
\iftag{prvTrue,defTrue}{$\ltrue$ या चिन्हाला वेगवेगळ्या ठिकाणी
  ``सत्यता'', ``वेरम'' किंवा ``शिखर'' म्हणतात.}{}
\end{intro}

\begin{defn}[सूत्र]
\ollabel{defn:formulas}
विधानीय तर्कशास्त्रातील \emph{!!{formula}s} चा संच~$\Frm[L_0]$
पुढीलप्रमाणे विगमनाने परिभाषित केला जातो:
\begin{enumerate}
\tagitem{prvFalse}{$\lfalse$ हे आण्विक !!{formula} आहे.}{}

\tagitem{prvTrue}{$\ltrue$ हे आण्विक !!{formula} आहे.}{}

\item प्रत्येक !!{propositional variable}~$\Obj p_i$ हे आण्विक
  !!{formula} आहे.

\tagitem{prvNot}{$!A$ हे !!a{formula} असेल, तर $\lnot !A$ हे
  !!a{formula} आहे.}{}

\tagitem{prvAnd}{$!A$ आणि $!B$ ही !!{formula}s असतील, तर $(!A \land
  !B)$ हे !!a{formula} आहे.}{}

\tagitem{prvOr}{$!A$ आणि $!B$ ही !!{formula}s असतील, तर $(!A \lor !B)$
  हे !!a{formula} आहे.}{}

\tagitem{prvIf}{$!A$ आणि $!B$ ही !!{formula}s असतील, तर $(!A \lif !B)$
  हे !!a{formula} आहे.}{}

\tagitem{prvIff}{$!A$ आणि $!B$ ही !!{formula}s असतील, तर $(!A \liff !B)$
  हे !!a{formula} आहे.}{}

\tagitem{limitClause}{इतर काहीही !!a{formula} नाही.}{}
\end{enumerate}
\end{defn}

\begin{explain}
!!{formula}s ची ही व्याख्या \emph{विगमनाधारित व्याख्या} आहे. मूलतः,
आपण !!{formula}s चा संच अनंतपणे अनेक टप्प्यांत तयार करतो. आरंभीच्या
टप्प्यात आपण सर्व आण्विक सूत्रांना सूत्रे घोषित करतो; हे व्याख्येतील
पहिल्या काही बाबींशी, म्हणजे
\iftag{prvTrue}{$\ltrue$, }{}%
\iftag{prvFalse}{$\lfalse$, }{}%
$\Obj p_i$ यांच्या बाबींशी जुळते. म्हणून ``आण्विक !!{formula}''
म्हणजे या रूपातील कोणतेही !!{formula}.

व्याख्येतील इतर बाबी आधीच तयार केलेल्या !!{formula}s पासून नवी
!!{formula}s तयार करण्याचे नियम देतात. दुसऱ्या टप्प्यात या नियमांनी
आण्विक !!{formula}s पासून !!{formula}s तयार करता येतात. तिसऱ्या
टप्प्यात आण्विक सूत्रे आणि दुसऱ्या टप्प्यात मिळालेली सूत्रे यांपासून
नवी सूत्रे तयार करतो, आणि ही प्रक्रिया पुढे चालू राहते. अशा एखाद्या
टप्प्यात शेवटी तयार होणारी प्रत्येक गोष्ट !!{formula} असते; इतर
काहीही सूत्र नसते.
\end{explain}

द्विस्थानी संयोजक~$\ast$ वापरून $!B$, $!C$ पासून तयार केलेले
सूत्र $(!B \ast !C)$ लिहिताना आपण अनेकदा सर्वांत बाहेरची कंसांची
जोडी वगळून केवळ~$!B \ast !C$ असे लिहू.

\begin{tagblock}{defNot,defOr,defAnd,defIf,defIff,defTrue,defFalse,defEx,defAll}
\begin{defn}
परिभाषित कारके वापरून तयार केलेली सूत्रे पुढीलप्रमाणे समजायची:

\begin{tagenumerate}{defTrue,defFalse,defNot,defOr,defAnd,defIf,defIff,defEx,defAll}

\tagitem{defTrue}{$\ltrue$ हा
  \iftag{prvFalse}{$\lnot\lfalse$}{$(!A \lor \lnot !A)$ (येथे~$!A$
    हे एखादे निश्चित आण्विक !!{formula} आहे)} चा संक्षेप आहे.}{}

\tagitem{defFalse}{$\lfalse$ हा
  \iftag{prvTrue}{$\lnot\ltrue$}{$(!A \land \lnot !A)$ (येथे~$!A$
    हे एखादे निश्चित आण्विक !!{formula} आहे)} चा संक्षेप आहे.}{}

\tagitem{defNot}{$\lnot !A$ हा $!A \lif \lfalse$ चा संक्षेप आहे.}{}

\tagitem{defOr}{$!A \lor !B$ हा
  \iftag{prvAnd}{$\lnot(\lnot !A \land \lnot !B)$}{$\lnot !A \lif
    !B$} चा संक्षेप आहे.}{}

\tagitem{defAnd}{$!A \land !B$ हा
  \iftag{prvOr}{$\lnot(\lnot !A \lor \lnot !B)$}{$\lnot (!A \lif
    \lnot !B)$} चा संक्षेप आहे.}{}

\tagitem{defIf}{$!A \lif !B$ हा
  \iftag{prvOr}{$\lnot !A \lor !B$}{$\lnot (!A \land \lnot !B)$} चा संक्षेप आहे.}{}
%READERNOTE{MRPL-001: गोठवलेल्या इंग्रजीतील पहिल्या संक्षेपात जुळणारा आरंभीचा कंस नसतानाही शेवटचा कंस आहे. मराठीत मानक नकार-विकल्पयोग सूत्र दिले असून इंग्रजी बाइट्स बदललेले नाहीत.}

\tagitem{defIff}{$!A \liff !B$ हा $(!A \lif !B) \land (!B
  \lif !A)$ चा संक्षेप आहे.}{}
\end{tagenumerate}
\end{defn}
\end{tagblock}

\begin{defn}[विन्यासात्मक अनन्यता]
$\ident$ हे चिन्ह चिन्हमालांमधील विन्यासात्मक अनन्यता व्यक्त करते;
म्हणजे $!A \ident !B$ तेव्हा आणि केवळ तेव्हाच, जेव्हा $!A$ आणि $!B$
या समान लांबीच्या चिन्हमाला असतात आणि प्रत्येक स्थानी त्यांच्यात तेच
चिन्ह असते.
\end{defn}

$\ident$ या चिन्हाच्या दोन्ही बाजूंना जोडणीने मिळालेल्या चिन्हमाला
येऊ शकतात. उदाहरणार्थ, $!A \ident (!B \lor !C)$ याचा अर्थ असा:
आरंभीचा कंस, चिन्हमाला~$!B$, $\lor$ हे चिन्ह, चिन्हमाला~$!C$ आणि
शेवटचा कंस या क्रमाने जोडून मिळणारी चिन्हमाला आणि चिन्हमाला~$!A$
या एकच आहेत. असे असेल, तर $!A$ चे पहिले चिन्ह आरंभीचा कंस आहे;
$!A$ मध्ये दुसऱ्या चिन्हापासून सुरू होणारी $!B$ ही उपचिन्हमाला आहे;
तिच्यानंतर $\lor$ येते; इत्यादी बाबी आपल्याला माहीत होतात.

\end{document}
